\documentclass[10pt,a4paper]{article}
\usepackage[margin=17mm]{geometry}
\usepackage{unicode-math}
\usepackage{experimental/finaldelim-draw}
\usepackage{fancyhdr}
\pagestyle{fancy}\fancyhf{}\fancyhead[L]{finaldelim: slashes and angles}\fancyfoot[C]{\thepage}
\setlength\parindent{0pt}
\newcommand\Cell[2]{\hbox to .105\textwidth{\hfil$\displaystyle\left#2\phantom{\vrule width1pt height\dimexpr#1pt/2+\Umathaxis\displaystyle\relax depth\dimexpr#1pt/2-\Umathaxis\displaystyle\relax}\right.$\hfil}}
\newcommand\Four[1]{\Cell{#1}{/}\Cell{#1}{\backslash}\Cell{#1}{\langle}\Cell{#1}{\rangle}}
\newcommand\Row[1]{\vbox to \dimexpr#1pt+15pt\relax{\vfil\hbox to \textwidth{\hbox to .12\textwidth{#1 pt\hfil}\FinalDelimDrawOff\Four{#1}\hfil\FinalDelimDrawOn\Four{#1}}\vfil}}
\newcommand\FontPage[1]{\clearpage\setmathfont{#1}\fancyhead[R]{#1}
{\Large\bfseries #1}\par\medskip
\hbox to \textwidth{\hskip.12\textwidth\hbox to .42\textwidth{\hfil Native glyphs\hfil}\hfil\hbox to .42\textwidth{\hfil Continuous outlines\hfil}}\par
\hbox to \textwidth{\hbox to .12\textwidth{Target\hfil}\Four{8}\hfil\Four{8}}\par\smallskip\hrule
\begingroup\FinalDelimSetup{q=0,scale=1,niji=false}\delimiterfactor=1000\delimitershortfall=0pt\nulldelimiterspace=0pt
\Row{10}\Row{14}\Row{18}\Row{24}\Row{32}\Row{48}\Row{64}\Row{96}
\endgroup}
\begin{document}
{\Huge\bfseries Slashes and angles}\par\bigskip
Pages 2--9: slashes and single angles. Pages 10--17: double angles.\par\medskip
Eight fonts; identical target heights for native glyphs and continuous outlines.
Each group shows slash, backslash, left angle and right angle.\par\bigskip
Inspect slant, width, stroke weight, the angle vertex, and behavior beyond the
largest prebuilt glyph. Requested height does not imply that a font can supply
that height: without an assembly, native output may stop growing. The new renderer
continues the recent width trend while retaining normal stroke thickness.\par\bigskip
Controlled sizing: q=0, scale=1, delimiterfactor=1000, delimitershortfall=0pt.
Beyond the prebuilt range, terminal and vertex contours are mapped around the extended stroke centerlines.\par\bigskip
\begin{tabular}{lr}Font & Page\\\hline
Latin Modern Math & 2\\
STIX Two Math & 3\\
TeX Gyre Pagella Math & 4\\
TeX Gyre Termes Math & 5\\
XITS Math & 6\\
Libertinus Math & 7\\
Fira Math & 8\\
GFS Neohellenic Math & 9\\
\end{tabular}
\FontPage{Latin Modern Math}
\FontPage{STIX Two Math}
\FontPage{TeX Gyre Pagella Math}
\FontPage{TeX Gyre Termes Math}
\FontPage{XITS Math}
\FontPage{Libertinus Math}
\FontPage{Fira Math}
\FontPage{GFS Neohellenic Math}

\newcommand\DCell[2]{\hbox to .20\textwidth{\hfil$\displaystyle\left#2\phantom{\vrule width1pt height\dimexpr#1pt/2+\Umathaxis\displaystyle\relax depth\dimexpr#1pt/2-\Umathaxis\displaystyle\relax}\right.$\hfil}}
\newcommand\DRow[1]{\vbox to \dimexpr#1pt+15pt\relax{\vfil\hbox to \textwidth{\hbox to .12\textwidth{#1 pt\hfil}\FinalDelimDrawOff\DCell{#1}{\lAngle}\DCell{#1}{\rAngle}\hfil\FinalDelimDrawOn\DCell{#1}{\lAngle}\DCell{#1}{\rAngle}}\vfil}}
\newcommand\DoublePage[1]{\clearpage\setmathfont{#1}\fancyhead[R]{#1}
{\Large\bfseries Double angles: #1}\par\medskip
Native left group; drawn right group. Each group: left and right double angle.
The two strokes extend independently; their source offsets are retained.\par\bigskip
\begingroup\FinalDelimSetup{q=0,scale=1,niji=false}\delimiterfactor=1000\delimitershortfall=0pt\nulldelimiterspace=0pt
\DRow{32}\DRow{64}\DRow{96}\DRow{144}\endgroup}
\DoublePage{Latin Modern Math}
\DoublePage{STIX Two Math}
\DoublePage{TeX Gyre Pagella Math}
\DoublePage{TeX Gyre Termes Math}
\DoublePage{XITS Math}
\DoublePage{Libertinus Math}
\DoublePage{Fira Math}
\DoublePage{GFS Neohellenic Math}
\end{document}
